%%% SECTION 10: PATIENT INSTRUCTIONS FOR PHYSICAL AI ONCOLOGY TRIALS %%%

[This section should be approximately 8-10 pages. The major theme is that
Patient Instructions: Physical AI Trials \cite{patient-instructions-paper}
provides clear, accessible, patient-facing documentation for participants in
Physical AI oncology trials. This document demonstrates that control has shifted
towards the patient's side, with more patient rights and comprehensive
information about every robot type they may encounter. The patient instructions
complement the informed consent requirements of Physical AI Adaption: 21 CFR
Part 50 \cite{cfr50-adapt} and the patient rights provisions of AB 2847 from
the Clinical Trial Site Complete Documentation \cite{site-docs}.]

[PRIMARY SOURCE: national-platform/patient\_robot/ directory containing two
chunked files: 01\_preamble\_robots1to5.tex (lines 1-197) and
02\_robots6to10\_closing.tex (lines 199-371). These contain the patient
instructions for ten robot types.]

\subsection{Purpose and Design Philosophy}
\label{subsec:instructions-purpose}

[Build from national-platform/patient\_robot/01\_preamble\_robots1to5.tex,
the preamble section. Explain the design philosophy behind the patient
instructions: plain language, clear formatting, consistent structure for each
robot type, emphasis on patient rights and comfort. The instructions were
designed to be accessible to patients of all education levels and to reduce
anxiety about interacting with Physical AI systems.]

[MAJOR THEME: Patient instructions are a critical tool for shifting control
towards patients. By providing comprehensive, understandable information
about each robot type, patients can make informed decisions about their
participation and exercise their rights effectively. Reference NCI information
on clinical trial safety \cite{nci-trials-safety} for context on patient
concerns about trial participation.]

\subsection{Surgical Robots}
\label{subsec:instructions-surgical}

[Build from national-platform/patient\_robot/01\_preamble\_robots1to5.tex,
Robot Type 1 (Surgical Robots). Cover what the patient needs to know about
surgical robots including: what the robot does, how it interacts with the
patient, what the patient will experience, safety features, and the patient's
rights during the procedure. Cross-reference with USL surgical robot scores
from Physical AI Unification Standard Level \cite{usl-standard}.]

\subsection{Collaborative Robots (Cobots)}
\label{subsec:instructions-cobots}

[Build from national-platform/patient\_robot/01\_preamble\_robots1to5.tex,
Robot Type 2 (Cobots). Cover cobot interactions including medication
preparation assistance, sample handling, and patient support tasks. Explain
the safety features specific to cobots working in close proximity to patients.]

\subsection{Radiotherapy Positioning Robots}
\label{subsec:instructions-rt-positioning}

[Build from national-platform/patient\_robot/01\_preamble\_robots1to5.tex,
Robot Type 3 (RT Positioning Robots). Cover how these robots assist with
precise patient positioning for radiation therapy, the importance of
accuracy, and what the patient will experience during positioning.]

\subsection{Needle-Placement Robots}
\label{subsec:instructions-needle}

[Build from national-platform/patient\_robot/01\_preamble\_robots1to5.tex,
Robot Type 4 (Needle-Placement Robots). Cover needle-guided procedures
including biopsies and targeted drug delivery, precision advantages, and
patient comfort measures.]

\subsection{Social Companion Robots}
\label{subsec:instructions-companion}

[Build from national-platform/patient\_robot/01\_preamble\_robots1to5.tex,
Robot Type 5 (Social Companion Robots). Cover the role of companion robots
in providing emotional support, information delivery, and monitoring patient
well-being during trial participation.]

[MAJOR THEME: Social companion robots represent a unique aspect of Physical AI
oncology trials that has no analog in traditional trials. These robots address
the emotional and psychological needs of cancer patients, demonstrating that
Physical AI is not just about clinical efficiency but also about patient
well-being and convenience.]

\subsection{Humanoid Robots}
\label{subsec:instructions-humanoid}

[Build from national-platform/patient\_robot/02\_robots6to10\_closing.tex,
Robot Type 6 (Humanoids). Cover the multi-functional role of humanoid robots
in trial sites including patient guidance, logistics support, and general
assistance. Explain the unique considerations for patient interactions with
human-shaped robots.]

\subsection{Radiotherapy Motion-Tracking Robots}
\label{subsec:instructions-rt-tracking}

[Build from national-platform/patient\_robot/02\_robots6to10\_closing.tex,
Robot Type 7 (RT Motion-Tracking Robots). Cover how these robots track
patient movement during radiation therapy to ensure treatment accuracy,
what the patient will experience, and safety interlocks.]

\subsection{Imaging Robots}
\label{subsec:instructions-imaging}

[Build from national-platform/patient\_robot/02\_robots6to10\_closing.tex,
Robot Type 8 (Imaging Robots). Cover robotic imaging systems for diagnostic
and monitoring purposes, including robotic CT, MRI assistance, and
ultrasound guidance. Reference DICOM \cite{dicom} as the imaging standard.]

\subsection{Steerable Needle Robots}
\label{subsec:instructions-steerable}

[Build from national-platform/patient\_robot/02\_robots6to10\_closing.tex,
Robot Type 9 (Steerable Needle Robots). Cover advanced needle guidance
systems for complex anatomical targets, precision advantages over manual
needle placement, and patient experience.]

\subsection{Rehabilitation Exoskeletons}
\label{subsec:instructions-rehab}

[Build from national-platform/patient\_robot/02\_robots6to10\_closing.tex,
Robot Type 10 (Rehab Exoskeletons). Cover the role of rehabilitation
exoskeletons in post-treatment recovery for oncology patients, including
mobility assistance, strength training, and recovery monitoring.]

\subsection{Patient Benefits Summary}
\label{subsec:instructions-benefits}

[Synthesize the patient benefits across all ten robot types. MAJOR THEME:
Physical AI oncology trials offer patients multiple benefits including:
convenience (reduced wait times, streamlined scheduling), lower cost of
treatments (reduced staffing overhead passed to patients as lower costs),
higher quality of treatment (precision robotics, continuous monitoring),
emotional support (companion robots), faster recovery (rehabilitation
exoskeletons), and comprehensive data-driven care. Reference NCI information
on trial costs \cite{nci-trials-paying} for comparison context.]

[Connect patient instructions to the broader National Platform by showing
how they implement the patient rights provisions of Physical AI Adaption:
21 CFR Part 50 \cite{cfr50-adapt}, the consent requirements of Adaption: ICH
Harmonised Guideline \cite{ich-adapt}, and the patient rights legislation
(AB 2847) from Clinical Trial Site Complete Documentation \cite{site-docs}.]
